\documentclass{article}%
\usepackage[T1]{fontenc}%
\usepackage[utf8]{inputenc}%
\usepackage{lmodern}%
\usepackage{textcomp}%
\usepackage{lastpage}%
\usepackage{geometry}%
\geometry{margin=2.2cm}%
%
\usepackage{amsmath}%
\usepackage{amssymb}%
\usepackage{booktabs}%
\usepackage{longtable}%
\usepackage{array}%
\usepackage{hyperref}%
\hypersetup{hidelinks}%
\title{IME 2027 {-} Física (claude\_haiku\_5\_5) x Poliedro\\[4pt]\large Comparação com o gabarito Poliedro\thanks{Gerado com gpt-resolve: \url{https://github.com/lgabs/gpt-resolve}}}%
\date{10/11/2026}%
%
\begin{document}%
\normalsize%
\maketitle%
\section*{Resumo}\textbf{Nota média: 10,0/10} (30 de 30 itens com a resolução correta; cada item vale o mesmo dentro da questão).\par\medskip%
\begin{longtable}{@{}c c p{0.78\textwidth}@{}}\toprule Questão & Nota & Resumo \\ \midrule
1 & 10,0 & A resolução está correta nos três itens e bate com o gabarito Poliedro: a) 250 m/s, b) 4000 K, c) 864 kJ. Refiz as contas e todas conferem. A resolução vai além do gabarito: verifica que o projeto original é coerente, mostra que o foguete ainda sobe na vertical do alvo (v\_y $\approx$ 64 m/s) e confere o item c pela relação de Mayer. Não há divergência. O único ponto a levar ao professor é a convenção de sinal do trabalho no item c, que é de redação e pouco importa. \\
2 & 10,0 & A resolução e o gabarito Poliedro coincidem nos três itens, e os dois estão certos. Refiz a questão resolvendo o sistema completo de equilíbrio dos nós (12$\times$12, frações exatas). a) R\_A = 200 kN para cima e R\_B = 100 kN para baixo, sem componentes horizontais. b) CD = 75 kN (compressão), CF = 125 kN (tração), DF = 200 kN (compressão), AD = 200 kN (compressão), AE = 0, DE = 75 kN (compressão), BE = 100 kN (tração), EF = 125 kN (tração). c) AE é barra de contraventamento: tem esforço nulo sob P, mas sem ela ADEB vira mecanismo (posto 11 < 12). A resolução também mostra que N\_AE = 5H/3 sob uma força horizontal H, o que conferi numericamente. Restam só imprecisões de redação, sem impacto na nota. \\
3 & 10,0 & A resolução e o gabarito do Poliedro coincidem nos três itens: a) R\_A = 300 N, vertical para cima (T$_{0}$ = 300 N); b) $\Delta$T = 100 $^\circ$C; c) Q = 5$\cdot$10$^{6}$ J, com rompimento em 37,5\% do dia. Refiz as contas em Python e confirmei os valores, e a física dos dois está correta (braço fixo L para a tração e L'/2 para o peso, integração linear da irradiância no ramo crescente). Não há divergência real. A única diferença é de comentário sobre a dilatação da área (0,2\% na área contra 0,05\% no tempo), e as duas afirmações são coerentes. Questão 3: 3/3 itens corretos. \\
4 & 10,0 & Resolução e gabarito Poliedro concordam nos três itens, e os dois estão certos. (a) U = 900 V, i1 = 50 A, i2 = 0. (b) A razão é 2,5 e o fusível rompe em $\approx$ 1,5 s; pela minha medição na figura, $\approx$ 1,52 s. (c) A corrente do gerador no curto é 500 A, a razão é 25 e o patamar dá 100 ms > 50 ms, então o fusível não protege a fonte. Os pontos em aberto são de importância baixa: a extrapolação do patamar além de x = 5, a precisão da leitura gráfica e a ressalva hipotética (curva em 1/x$^2$) que a resolução acrescenta no item (c). \\
5 & 10,0 & Os três itens da questão 5 batem com o gabarito Poliedro, e as contas conferem em Python. a) T = 450 K. b) Q = 90 kW. c) Seriam necessários 116,7 kW; o aquecimento consome 33,3 kW; sobram no máximo 71,7 kW, o que atende só 215 residências; $\Delta$$\dot S$ $\approx$ +0,044 kW/K, então a 2ª lei não é violada. A conclusão é a mesma: a promessa das 350 residências não é realista. A resolução vai além do gabarito com limites corretos (85 kW no caso reversível, 100 kW de teto absoluto). Também ressalva que a água precisaria entrar abaixo de 290 K, o que é um refinamento válido e não contradiz o gabarito. Não há divergência de mérito. \\
6 & 10,0 & A questão 6 está totalmente de acordo com o gabarito Poliedro nos 4 itens: v = 68 m/s; R = 5$\surd$3 $\approx$ 8,66 m; T $\approx$ 1,07$\times$10$^3$ N; V $\approx$ 55$\times$10$^{-}$$^3$ m$^3$ (E $\approx$ 551 N). Refiz as contas em Python e elas confirmam os valores. A resolução chega a R pela distância |BO| = 20 m, enquanto o gabarito usa o ângulo de 60$^\circ$. É a mesma relação. Não achei erro de física nem de conta em nenhum dos dois. Ficam apenas hipóteses razoáveis e não declaradas no texto: o fio preso na vertical do centro e o som chegando a O pelo ar. \\
7 & 10,0 & A resolução está certa nos três itens e concorda com o gabarito Poliedro: (a) 0,45 m; (b) |razão| $\approx$ 1,6 (29/18), com kx $\approx$ 6,44 N; (c) $\surd$2 m/s. Refiz as contas e elas conferem. A única divergência é o sinal no item b: a resolução dá a razão com |$\Delta$U\_g| (+1,6), embora tenha escrito $\Delta$U\_g = -2x\_s, e o gabarito usa a variação com sinal (-1,6), que é a leitura mais literal do enunciado. Se a banca exigir o sinal, esse ponto pode custar nota. \\
8 & 10,0 & A resolução está correta nos quatro itens e concorda com o gabarito Poliedro. a) O anteparo se afasta da rede. b) D(t1) = 0,2$\surd$7 $\approx$ 0,53 m. c) v $\approx$ 0,255 m/s com raízes exatas, contra 0,254 m/s do gabarito, que usa a tabela. d) $\Delta$s $\approx$ 1,79 m, contra 1,792 m do gabarito. As únicas diferenças vêm do arredondamento das raízes da tabela. Conferi o enunciado transcrito com o PDF original e ele está fiel. Fica só um pequeno deslize de redação na resolução: diz \textbackslash{}"centésimos de m/s\textbackslash{}" quando o valor está na casa dos décimos. \\
9 & 10,0 & Questão 9: a resolução está correta e bate com o gabarito do Poliedro nos dois itens. Item a: u = c($\surd$3h - L)/($\surd$3L - h). Item b: L/$\surd$3 < h < L. O enunciado do PDF (página 9) confere com a transcrição. Refiz a geometria da reflexão (i = r = 30$^\circ$, normal a 30$^\circ$ da vertical) e chego ao mesmo sistema linear do gabarito, cuja solução é ct = $\surd$3L - h e ut = $\surd$3h - L. Conferi numericamente com quatro pares (L, h) e o resultado se confirma. A resolução é mais completa que o gabarito: confirma que o impacto é o primeiro contato e descarta o outro sinal da normal. Sobra um único ponto para o professor, de baixa importância: se a condição do item b deveria incluir u < ($\surd$3/2)c. O enunciado manda desprezar o movimento da superfície, então a resposta esperada não muda. \\
10 & 10,0 & A resolução confere com o gabarito Poliedro nos dois itens, e os dois estão corretos. No item (a), só a Situação 2 serve: as Situações 1 e 4 têm campo diferente de zero no eixo e a 3 é instável. No item (b), p = ($\pi$d/2)$\cdot$$\surd$($\pi$mv/($\mu$0 i q)), que é meio período do MHS horizontal vezes v. Refiz a dedução e uma integração numérica com o campo exato, e o resultado da fórmula se confirmou com desvio de 3$\cdot$10$^{-}$$^{4}$, que é exatamente a correção de segunda ordem prevista. A resolução justifica cada situação, como o item (a) exige. Não encontrei erro. \\
\bottomrule\end{longtable}%
\section*{Principais pontos de dúvida}%
\begin{itemize}%
\item%
\textbf{Q7} \textit{(solution)} Item b, sinal da razão: o enunciado pede a razão entre E\_el e a 'variação da energia potencial gravitacional de P2'. Essa variação é negativa, porque P2 desce. A resolução responde E\_el/|$\Delta$U\_g| $\approx$ +1,6, embora tenha escrito $\Delta$U\_g = -2x\_s. O gabarito Poliedro dá -1,6. Se a banca exigir o sinal, o item pode perder nota. Convém ver com um professor se a resposta em módulo é aceita.%
\end{itemize}%
\section*{Questão 1}%
\href{https://poliedroresolve.sistemapoliedro.com.br/vestibulares/ime/2027/ime-2-fase-dia-2-27-10-2026-2-fase-dia-2/dissertativa-1-fisica-2-fase-dia-2-ime-2027}{Resolução do gabarito}\par\smallskip%
\begin{longtable}{@{}p{0.06\textwidth} p{0.33\textwidth} p{0.33\textwidth} p{0.1\textwidth} p{0.08\textwidth}@{}}\toprule Item & Resolução & Gabarito & Veredito & Certo \\ \midrule
a & v = 250 m/s & v0 = 250 m/s & igual & ambos \\
b & T = 4000 K & T\_falha = 4000 K & igual & ambos \\
c & W = 864 kJ (trabalho realizado pelos gases, positivo) & W = 864 kJ (Q = 2880 kJ, $\Delta$U = 2016 kJ) & igual & ambos \\
\bottomrule\end{longtable}%
\begin{itemize}%
\item%
\textbf{a)} Com $\theta$ = 45$^\circ$, a equação da trajetória fica y = x - g$\cdot$x$^2$/v$^2$. Substituindo x = 2000 m e y = 1360 m, vem 4$\cdot$10$^{7}$/v$^2$ = 640, logo v$^2$ = 62500 e v = 250 m/s (conferido em Python). A resolução ainda mostra que o foguete continua subindo: v\_y$^2$ = 4050 > 0 e o ápice fica em x = 3125 m, depois do alvo. Também confirma que o projeto original é coerente (H = 1000 m e x\_V = 2000 m).%
\item%
\textbf{b)} T/T0 = (v/v0)$^2$, então T = 2560$\cdot$(250/200)$^2$ = 2560$\cdot$25/16 = 4000 K. A resolução observa, com razão, que a proporcionalidade vale na escala absoluta.%
\item%
\textbf{c)} Q = m$\cdot$cp$\cdot$$\Delta$T = 2,0$\cdot$1,0$\cdot$1440 = 2880 kJ. Como W = Q - $\Delta$U, W = 2880 - 2016 = 864 kJ. A resolução confere pela relação de Mayer: cv = 0,70 kJ/(kg$\cdot$K), e m$\cdot$(cp - cv)$\cdot$$\Delta$T = 864 kJ. Também verifica que $\gamma$ $\approx$ 1,43 é um valor plausível. A convenção de sinal (trabalho feito pelo gás, positivo na expansão) é a mesma do gabarito.%
\end{itemize}%
\paragraph{Pontos de dúvida}%
\begin{itemize}%
\item%
\textit{(statement)} O item c pede 'o trabalho realizado durante o aquecimento' sem dizer se é o trabalho feito pelo gás ou sobre ele. A resolução e o gabarito adotam o trabalho feito pelo gás (+864 kJ, expansão isobárica). Com a convenção contrária, a resposta seria -864 kJ.%
\item%
\textit{(solution)} A hipótese de que o foguete se comporta como projétil logo ao sair do lançador (sem empuxo depois) fica implícita. É o que o enunciado pede ao fornecer a velocidade de lançamento e mandar desprezar a resistência do ar, e o gabarito usa a mesma hipótese.%
\end{itemize}%
\section*{Questão 2}%
\href{https://poliedroresolve.sistemapoliedro.com.br/vestibulares/ime/2027/ime-2-fase-dia-2-27-10-2026-2-fase-dia-2/dissertativa-2-fisica-2-fase-dia-2-ime-2027}{Resolução do gabarito}\par\smallskip%
\begin{longtable}{@{}p{0.06\textwidth} p{0.33\textwidth} p{0.33\textwidth} p{0.1\textwidth} p{0.08\textwidth}@{}}\toprule Item & Resolução & Gabarito & Veredito & Certo \\ \midrule
a & R\_A = 200 kN, vertical, para cima; R\_B = 100 kN, vertical, para baixo (B precisa ser ancorado); componentes horizontais R\_Ax = R\_Bx = 0. & A = 200 kN vertical para cima; B = 100 kN vertical para baixo (A\_x = B\_x = 0). & igual & ambos \\
b & N\_CD = 75 kN (C); N\_CF = 125 kN (T); N\_DF = 200 kN (C); N\_AD = 200 kN (C); N\_AE = 0; N\_DE = 75 kN (C); N\_BE = 100 kN (T); N\_EF = 125 kN (T). & N\_CD = 75 kN (C); N\_CF = 125 kN (T); N\_DF = 200 kN (C); N\_AD = 200 kN (C); N\_AE = 0; N\_DE = 75 kN (C); N\_BE = 100 kN (T); N\_EF = 125 kN (T). & igual & ambos \\
c & AE é barra de contraventamento. Sob P o esforço nela é nulo, mas sem ela o quadrilátero ADEB vira paralelogramo articulado (mecanismo horizontal, estrutura hipostática com m+r = 11 < 12). Com AE a estrutura fica rígida e isostática, e a barra trabalha com N\_AE = 5H/3 quando há força horizontal H. & N\_AE = 0 só para este carregamento. AE é necessária para outras distribuições de carga e garante a rigidez: sem ela ADEB vira paralelogramo articulado e o guindaste balança lateralmente. & igual & ambos \\
\bottomrule\end{longtable}%
\begin{itemize}%
\item%
\textbf{a)} Refiz a conta. Momentos em A: (-3)(-100) + 3$\cdot$B\_y = 0 $\Rightarrow$ B\_y = -100 kN; $\Sigma$F\_y $\Rightarrow$ A\_y = 200 kN. As reações horizontais são nulas. O gabarito justifica isso pelo fato de a única barra em B (BE) ser vertical. A resolução usa um corte em CDEF, em que só AE tem componente horizontal, logo N\_AE = 0 e R\_Ax = 0. As duas justificativas valem. Também resolvi o sistema completo 12$\times$12 com frações exatas: A\_x = 0, A\_y = 200, B\_x = 0, B\_y = -100.%
\item%
\textbf{b)} Todos os valores e naturezas coincidem. Conferi pelo sistema linear completo dos 6 nós (posto 12): CD = -75, CF = 125, DF = -200, AD = -200, AE = 0, DE = -75, BE = 100, EF = 125, com sinal positivo para tração. A resolução tira N\_AD e N\_BE do corte CDEF e confere nos nós D, E e A. O gabarito usa a sequência de nós C, F, D, E. Os dois caminhos são corretos.%
\item%
\textbf{c)} A conclusão é a mesma: a barra tem esforço nulo neste carregamento, mas é indispensável para a rigidez e o contraventamento lateral. Conferi numericamente: sem AE, a matriz de equilíbrio tem posto 11 para 12 equações, ou seja, é um mecanismo. Com uma força horizontal unitária em E para a direita, N\_AE = +5/3, o que confirma N\_AE = 5H/3 em tração. A resolução vai além do gabarito ao quantificar esse esforço, sem contradizê{-}lo.%
\end{itemize}%
\paragraph{Pontos de dúvida}%
\begin{itemize}%
\item%
\textit{(solution)} No item c, a resolução fala em 'DE e AB (horizontais, 3 m)' como lados do paralelogramo articulado. AB não é barra: é o trecho de solo entre os dois apoios fixos. O argumento continua certo (quadrilátero articulado de quatro barras), mas a redação é imprecisa.%
\item%
\textit{(statement)} Os dois lados supõem apoios fixos (articulados, com duas reações cada) em A e B, como sugerem os triângulos idênticos da figura. Se B fosse um apoio móvel, B\_x = 0 seria imposto de saída e a contagem de estaticidade mudaria. O resultado sob P não muda, porque B\_x já é nulo.%
\item%
\textit{(solution)} R\_B aponta para baixo, ou seja, o apoio B é tracionado (arrancamento). A resolução observa isso e diz que B deve ser ancorado. O gabarito não comenta. A informação é correta e só complementa a resposta.%
\end{itemize}%
\section*{Questão 3}%
\href{https://poliedroresolve.sistemapoliedro.com.br/vestibulares/ime/2027/ime-2-fase-dia-2-27-10-2026-2-fase-dia-2/dissertativa-3-fisica-2-fase-dia-2-ime-2027}{Resolução do gabarito}\par\smallskip%
\begin{longtable}{@{}p{0.06\textwidth} p{0.33\textwidth} p{0.33\textwidth} p{0.1\textwidth} p{0.08\textwidth}@{}}\toprule Item & Resolução & Gabarito & Veredito & Certo \\ \midrule
a & N = 300 N, vertical, para cima, sem componente horizontal (T$_{0}$ = W/2 = 300 N) & T = 300 N; R\_A = 300 N vertical para cima & igual & ambos \\
b & $\Delta$T = 100 $^\circ$C (100 K), de T = T$_{0}$(1+$\alpha$$\Delta$T) = 1,001$\cdot$T$_{0}$ & $\Delta$T = 100 $^\circ$C & igual & ambos \\
c & Q = 5$\cdot$10$^{6}$ J; rompimento em t = 3$\cdot$10$^{4}$ s = 37,5\% do dia (ramo crescente) & Q = 5$\cdot$10$^{6}$ J; rompe às 37,5\% do dia (12,5\% após 25\%) & igual & ambos \\
\bottomrule\end{longtable}%
\begin{itemize}%
\item%
\textbf{a)} m = 2000$\cdot$(2$\cdot$0,25$\cdot$0,05) = 50 kg; P = 50$\cdot$12 = 600 N. Momento no pino: T$\cdot$L = P$\cdot$L/2 $\Rightarrow$ T = 300 N; R\_A = P - T = 300 N, para cima. Refiz as contas e os dois chegam ao mesmo valor. A resolução ainda mostra que a componente horizontal da reação é nula, o que está certo, já que não há atrito na argola.%
\item%
\textbf{b)} O pino e o ponto de sustentação (argola) ficam fixos, então o braço da tração continua L e o do peso passa a L'/2. Daí T$\cdot$L = P$\cdot$L(1+$\alpha$$\Delta$T)/2 $\Rightarrow$ T = T$_{0}$(1+$\alpha$$\Delta$T). Com T\_max = 1,001$\cdot$T$_{0}$: $\alpha$$\Delta$T = 10$^{-}$$^3$ $\Rightarrow$ $\Delta$T = 100 $^\circ$C. O raciocínio é o mesmo nos dois.%
\item%
\textbf{c)} Q = 50$\cdot$1000$\cdot$100 = 5$\cdot$10$^{6}$ J. No ramo crescente, I = 0,2(t - t$_{0}$) W/m$^2$ com t em s, e A = 0,5 m$^2$, o que dá E = 0,05(t - t$_{0}$)$^2$. Igualando a Q: t - t$_{0}$ = 10$^{4}$ s, ou seja, 12,5\% do dia, e o horário é 25\% + 12,5\% = 37,5\%. O Poliedro faz a conta em \% do dia (E = 32000$\cdot$t$^2$ $\Rightarrow$ t = 12,5) e chega ao mesmo resultado. A resolução também confere que até o pico a energia absorvida seria 2$\cdot$10$^{7}$ J, maior que Q, o que confirma o ramo crescente. Sobre a dilatação da área: os 0,2\% citados pela resolução são a variação final da área, e os 0,05\% do Poliedro são o efeito disso no tempo de rompimento (cerca de 0,1\% a mais de energia, logo cerca de 0,05\% no tempo). As duas observações são coerentes e não mudam a resposta.%
\end{itemize}%
\paragraph{Pontos de dúvida}%
\begin{itemize}%
\item%
\textit{(statement)} A frase 'só conseguirá suportar 0,1\% a mais de tração' foi lida como T\_max = 1,001$\cdot$T$_{0}$. Essa é a leitura natural e é a que o gabarito usa.%
\item%
\textit{(statement)} O braço da tração fica igual a L porque a observação diz que o ponto de sustentação fica fixo e a haste continua vertical. Quem tomasse a tração aplicada na extremidade do objeto, que se desloca para L', teria T constante e nenhum rompimento. O enunciado afasta essa leitura, e tanto a resolução quanto o gabarito usam braço L.%
\item%
\textit{(statement)} Do gráfico se leem o pico de 4000 W/m$^2$ em 50\% e a variação linear de 25\% a 50\%. A conta também supõe incidência normal na face superior, toda a energia indo para o aquecimento (sem perdas) e temperatura uniforme no objeto. São hipóteses implícitas, as mesmas no gabarito e na resolução.%
\item%
\textit{(key)} Na resolução do Poliedro a dilatação superficial 'mudaria só 0,05\%'. Esse número é o efeito no tempo de rompimento, não a variação da área, que chega a 2$\alpha$$\Delta$T = 0,2\% no fim. A redação pode confundir, mas a conclusão está certa: o efeito é desprezível.%
\end{itemize}%
\section*{Questão 4}%
\href{https://poliedroresolve.sistemapoliedro.com.br/vestibulares/ime/2027/ime-2-fase-dia-2-27-10-2026-2-fase-dia-2/dissertativa-4-fisica-2-fase-dia-2-ime-2027}{Resolução do gabarito}\par\smallskip%
\begin{longtable}{@{}p{0.06\textwidth} p{0.33\textwidth} p{0.33\textwidth} p{0.1\textwidth} p{0.08\textwidth}@{}}\toprule Item & Resolução & Gabarito & Veredito & Certo \\ \midrule
a & U = 900 V; I\_Rec1 = 50 A (entrando pelo polo positivo, consumindo); I\_Rec2 = 0; corrente no gerador/fusível I = 50 A & U\_AB = 900 V; i1 = 50 A; i2 = 0 (tensão A--B igual à fcem de Rec2) & igual & ambos \\
b & Sim, rompe: I/I\_F = 50/20 = 2,5 $\Rightarrow$ t $\approx$ 1,5 s (leitura entre 1,45 e 1,55 s) & I/I\_fus = 2,5 $\Rightarrow$ fusível rompe em t $\approx$ 1,5 s (lido na curva) & igual & ambos \\
c & Não protege: U = 0, I = 1000/2 = 500 A, razão 25, patamar $\Rightarrow$ t\_rup = 0,1 s = 100 ms > 50 ms (sob a hipótese de patamar para x > 3,3) & I\_curto = 500 A, I/I\_fus = 25 $\Rightarrow$ t = 100 ms > 50 ms $\Rightarrow$ o fusível NÃO protege a fonte & igual & ambos \\
\bottomrule\end{longtable}%
\begin{itemize}%
\item%
\textbf{a)} Refeito: (1000-U)/2 = (U-850) + (U-900)/2 $\Rightarrow$ U = 900 V $\Rightarrow$ I = 50 A, i1 = 50 A, i2 = 0. Pela Figura 1, o polo positivo das três fontes fica do lado de A, o que confirma a convenção de sinais. O balanço de potências da resolução (50 kW = 5 + 42,5 + 2,5 kW) também fecha.%
\item%
\textbf{b)} Medi a Figura 2 em pixels no PDF renderizado a 300 dpi. A curva tem uma quina em (1,10; 3,84 s) e chega ao patamar em (3,36; 0,1 s). Interpolando em x = 2,5, dá t $\approx$ 1,52 s, o que confirma $\approx$ 1,5 s. A grade da figura tem passo de 0,5 nos dois eixos, como a resolução diz. A justificativa também está correta: a corrente é constante (circuito puramente resistivo) e maior que I\_FUSÍVEL.%
\item%
\textbf{c)} Com A--B em curto, U = 0 e a corrente do gerador, que é a que passa pelo fusível, vale E/r = 500 A. A razão é 25, além do fim do eixo (x = 5). Gabarito e resolução estendem o patamar de 0,1 s e concluem que 100 ms > 50 ms, então o fusível não protege a fonte. As correntes -850 A e -450 A nos receptores, e os 1800 A no fio do curto, não passam pelo fusível e não mudam a conclusão. A resolução acrescenta uma ressalva hipotética (uma curva em 1/x$^2$ daria $\approx$ 15 ms), mas mantém a resposta final "não".%
\end{itemize}%
\paragraph{Pontos de dúvida}%
\begin{itemize}%
\item%
\textit{(statement)} No item (c), a razão I/I\_FUSÍVEL = 25 fica bem além do fim do eixo da Figura 2 (x = 5). Gabarito e resolução estendem o patamar de 0,1 s. É a leitura que o enunciado sugere, mas é uma extrapolação que não está declarada no enunciado.%
\item%
\textit{(solution)} A resolução acrescenta um cenário hipotético: se a curva caísse com 1/x$^2$ (I$^2$t constante), daria $\approx$ 15 ms e o fusível protegeria. Isso não muda a resposta final "não", mas um corretor pode ler como hesitação. Vale manter a ressalva curta e deixar claro que a resposta segue a curva dada.%
\item%
\textit{(key)} O valor t $\approx$ 1,5 s do item (b) vem de leitura gráfica. Medindo em pixels, obtive $\approx$ 1,52 s. Uma faixa como 1,4--1,6 s deve ser aceita por qualquer correção razoável.%
\item%
\textit{(solution)} No curto, os receptores aparecem como fontes ideais que fornecem -850 A e -450 A. Num motor real, a fcem cairia com a rotação. Isso não afeta a corrente no fusível, que é E/r = 500 A, nem a conclusão.%
\end{itemize}%
\section*{Questão 5}%
\href{https://poliedroresolve.sistemapoliedro.com.br/vestibulares/ime/2027/ime-2-fase-dia-2-27-10-2026-2-fase-dia-2/dissertativa-5-fisica-2-fase-dia-2-ime-2027}{Resolução do gabarito}\par\smallskip%
\begin{longtable}{@{}p{0.06\textwidth} p{0.33\textwidth} p{0.33\textwidth} p{0.1\textwidth} p{0.08\textwidth}@{}}\toprule Item & Resolução & Gabarito & Veredito & Certo \\ \midrule
a & T\_fA = 450 K ($\eta$\_A = 60/150 = 0,40 = 0,80$\cdot$$\eta$\_C, logo $\eta$\_C = 0,50) & $\eta$1 = 40\%; T = 450 K & igual & ambos \\
b & Q\_fA = 90 kW & Q rejeitado = 90 kW & igual & ambos \\
c & Não realistas. Para 350 residências seriam precisos 116,7 kW. Com o aquecimento (Q\_água = 33,3 kW a 300 K), $\dot W$\_B $\leq$ 71,7 kW, o que dá no máximo 215 residências. $\sigma$ = 2/45 kW/K $\approx$ 44 W/K > 0, então a 2ª lei não proíbe o aquecimento, mas a água teria de entrar abaixo de 290 K. Limites extras: 85 kW no reversível com a repartição dada (255 residências) e 100 kW como limite absoluto (300 residências). & P necessária = 116,7 kW; Q\_água = 33,3 kW; a máquina 2 fornece 71,7 kW < 116,7 kW, então NÃO supre as 350 residências (máx. 215); $\Delta$$\dot S$\_universo = +0,04 kW/K, logo não viola a 2ª lei & equivalente & ambos \\
\bottomrule\end{longtable}%
\begin{itemize}%
\item%
\textbf{a)} $\dot W$ = 80$\cdot$3/4 = 60 kW; $\eta$ = 60/150 = 0,4 = 0,8(1 - T/900), logo T = 450 K. Conferi com Python: está certo.%
\item%
\textbf{b)} Pela 1ª lei, 150 - 60 = 90 kW. A resolução ainda confere que uma máquina reversível rejeitaria 75 kW e que $\sigma$ = 1/30 kW/K > 0. Tudo coerente.%
\item%
\textbf{c)} Os números são todos iguais. P = 350$\cdot$240/720 = 350/3 $\approx$ 116,7 kW. Q\_água = (50/60)$\cdot$4000$\cdot$10 = 100/3 $\approx$ 33,3 kW. $\dot W$ = 150 - 45 - 33,3 = 215/3 $\approx$ 71,7 kW, que dá 215 residências. $\Delta$$\dot S$ = -1/6 + 1/10 + 1/9 = 2/45 $\approx$ 0,044 kW/K, que o gabarito arredonda para 0,04. A conclusão também coincide: as 350 residências são inviáveis e o aquecimento não viola a 2ª lei. Há uma diferença de forma: o gabarito toma Q\_300 igual ao calor da água, e a resolução toma Q\_300 $\geq$ esse valor, o que dá $\dot W$ $\leq$ 71,7 kW; o resultado é o mesmo. A resolução acrescenta limites corretos que o gabarito não tem: 85 kW no caso reversível (Q\_300 $\geq$ 20 kW) e 100 kW de teto absoluto. Também acrescenta a ressalva de que a água precisa entrar abaixo de 290 K. Isso é um refinamento físico válido, e não contradiz o gabarito.%
\end{itemize}%
\paragraph{Pontos de dúvida}%
\begin{itemize}%
\item%
\textit{(solution)} No aquecimento de água, a resolução conclui que ele 'não é garantido' e exige T\_e < 290 K, um dado que o enunciado não fornece. O gabarito trata o aquecimento simplesmente como compatível com a 2ª lei ($\Delta$$\dot S$ > 0). A física da ressalva está certa, e a resolução diz que a 2ª lei não proíbe o aquecimento. Ainda assim, um corretor pode esperar a afirmação direta de que o aquecimento é viável e só a potência é insuficiente.%
\item%
\textit{(statement)} O enunciado não diz se o calor rejeitado a 300 K é exatamente o calor usado na água ou se pode ser maior. O gabarito assume que é igual, e a resolução usa Q\_300 $\geq$ Q\_água. Nos dois casos o máximo é 71,7 kW, ou 215 residências.%
\item%
\textit{(key)} O gabarito escreve $\Delta$$\dot S$\_universo = +0,04 kW/K, que é um arredondamento de 2/45 $\approx$ 0,0444 kW/K ($\approx$ 44 W/K). É só arredondamento, não é erro.%
\item%
\textit{(solution)} A resolução calcula um 'limite absoluto' de 100 kW em que a repartição de calor entre as fontes frias fica livre. Isso relaxa a condição do enunciado de que a fonte T recebe Q/2 = 45 kW. É uma informação extra e correta como cota, mas não é necessária. O mesmo vale para a hipótese de que o consumo médio é cota inferior da demanda de pico.%
\end{itemize}%
\section*{Questão 6}%
\href{https://poliedroresolve.sistemapoliedro.com.br/vestibulares/ime/2027/ime-2-fase-dia-2-27-10-2026-2-fase-dia-2/dissertativa-6-fisica-2-fase-dia-2-ime-2027}{Resolução do gabarito}\par\smallskip%
\begin{longtable}{@{}p{0.06\textwidth} p{0.33\textwidth} p{0.33\textwidth} p{0.1\textwidth} p{0.08\textwidth}@{}}\toprule Item & Resolução & Gabarito & Veredito & Certo \\ \midrule
a & v = 68 m/s & v = 68 m/s & igual & ambos \\
b & r = 5$\surd$3 $\approx$ 8,66 m & $\theta$ = 60$^\circ$ (cos$\theta$ = 1/2); R = 5$\surd$3 m & igual & ambos \\
c & T $\approx$ 1,07$\times$10$^3$ N (1,6$\cdot$667,4 $\approx$ 1068 N) & $\alpha$ = 60$^\circ$; T $\cong$ 1,07$\cdot$10$^3$ N & igual & ambos \\
d & E $\approx$ 551 N; V\_sub $\approx$ 5,5$\times$10$^{-}$$^2$ m$^3$ ($\approx$ 55 L) & E $\approx$ 551 N; V $\cong$ 55$\cdot$10$^{-}$$^3$ m$^3$ & igual & ambos \\
\bottomrule\end{longtable}%
\begin{itemize}%
\item%
\textbf{a)} Em A a velocidade, no sentido anti{-}horário, é tangente e aponta para O. Então f0$\cdot$c/(c-v) = 5f0/4, o que dá v = c/5 = 68 m/s. Conferido numericamente: 340/272 = 1,25.%
\item%
\textbf{b)} A resolução calcula D = |BO| = 10vd/c = 20 m e depois r$^2$ = (D$^2$ - d$^2$)/4 = 75. O gabarito usa o ângulo $\theta$ entre a velocidade em B e a reta BO: v$\cdot$cos$\theta$ = 34, então $\theta$ = 60$^\circ$, e tg$\theta$ = 2R/d dá R = 5$\surd$3. É a mesma conta, já que cos$\theta$ = d/D = 1/2. Conferido: f\_B/f0 = 340/374 = 10/11.%
\item%
\textbf{c)} Com sen$\alpha$ = r/L = $\surd$3/2 e T$\cdot$sen$\alpha$ = mv$^2$/r = 4624/5 = 924,8 N, sai T = (8/5)$\cdot$34$^2$/$\surd$3 $\approx$ 1067,9 N $\approx$ 1,07$\times$10$^3$ N. A resolução ainda justifica que o fio é preso no fundo, na vertical do centro, porque a tração é a única força horizontal e precisa ser centrípeta. Isso confere com a linha tracejada da vista lateral.%
\item%
\textbf{d)} E = mg + T$\cdot$cos$\alpha$ $\approx$ 17,3 + 533,9 $\approx$ 551,25 N, e V = E/($\rho$g) $\approx$ 0,0551 m$^3$. As duas formas de escrever (5,5$\times$10$^{-}$$^2$ m$^3$ e 55$\cdot$10$^{-}$$^3$ m$^3$) são o mesmo número.%
\end{itemize}%
\paragraph{Pontos de dúvida}%
\begin{itemize}%
\item%
\textit{(statement)} O texto não diz em que ponto do fundo o fio está preso. Tanto a resolução quanto o gabarito supõem que é na vertical do centro da trajetória. A figura (linha tracejada vertical na vista lateral) e a dinâmica (a única força horizontal precisa ser centrípeta) sustentam essa hipótese.%
\item%
\textit{(statement)} Ambos supõem que o som chega a O pelo ar (v\_som = 340 m/s) e que vale a posição da fonte no instante da emissão em A e B ("quando o objeto emite ao passar"). Essa leitura é coerente com o enunciado, que só dá a velocidade do som no ar.%
\item%
\textit{(solution)} A resolução acrescenta a hipótese f0 ? v/(2$\pi$r) $\approx$ 1,25 Hz para usar o Doppler instantâneo e conclui que a profundidade da piscina é H = 5 m. São complementos corretos e não alteram as respostas.%
\end{itemize}%
\section*{Questão 7}%
\href{https://poliedroresolve.sistemapoliedro.com.br/vestibulares/ime/2027/ime-2-fase-dia-2-27-10-2026-2-fase-dia-2/dissertativa-7-fisica-2-fase-dia-2-ime-2027}{Resolução do gabarito}\par\smallskip%
\begin{longtable}{@{}p{0.06\textwidth} p{0.33\textwidth} p{0.33\textwidth} p{0.1\textwidth} p{0.08\textwidth}@{}}\toprule Item & Resolução & Gabarito & Veredito & Certo \\ \midrule
a & d\_min = 0,45 m (retorno em x = 0,55 m a partir de B) & d = 0,45 m & igual & ambos \\
b & E\_el/|$\Delta$U\_g(P2)| $\approx$ 1,6 (= 29/18), com kx\_s = 58/9 $\approx$ 6,44 N e $\Delta$U\_g = -2x\_s & kx $\cong$ 6,44 N; U\_el/$\Delta$U\_g $\cong$ -1,6 & equivalente & ambos \\
c & v2 = $\surd$2 m/s $\approx$ 1,41 m/s, horizontal, de B para C, quando P2 atinge C & v2 = $\surd$2 m/s & igual & ambos \\
\bottomrule\end{longtable}%
\begin{itemize}%
\item%
\textbf{a)} Refiz a conta: com P2 fixa, a energia de P1 se conserva e dá 2(1-d) = 0,9(1/d - 1), ou seja, (1-d)(2 - 0,9/d) = 0. Daí d = 0,45 m. A resolução também separa corretamente o ponto de força nula (r = $\surd$0,45 $\approx$ 0,67 m, onde a velocidade é máxima) do ponto de retorno.%
\item%
\textbf{b)} O módulo bate: kx = 0,9/0,45$^2$ + 2 = 58/9 $\approx$ 6,44 N, E\_el = (kx)$\cdot$x/2 e $\Delta$U\_g = -mg$\cdot$x$\cdot$sen30$^\circ$ = -2x. Logo E\_el/$\Delta$U\_g = -(58/9)/4 = -29/18 $\approx$ -1,61. A resolução escreve $\Delta$U\_g = -2x\_s explicitamente, mas dá a resposta como razão com |$\Delta$U\_g|, por isso fica +1,6. O gabarito usa a variação com sinal e fica -1,6. A física e a conta são as mesmas; só muda a convenção de sinal. Ao pé da letra, 'variação da energia potencial gravitacional' tem sinal, então a forma do gabarito (-1,6) é a mais literal.%
\item%
\textbf{c)} Refiz a conta: com BC horizontal, o momento se conserva e as massas são iguais, então os deslocamentos são iguais e opostos. P2 anda 0,15 m até C, P1 anda 0,15 m para B e a separação vai de 0,45 m para 0,75 m. Pela energia: 0,9(1/0,45 - 1/0,75) = 0,8 J = 0,4v$^2$, logo v$^2$ = 2 e v = $\surd$2 m/s. Também confirmei que C é o primeiro vértice atingido, porque P1 ainda fica a 0,40 m de B.%
\end{itemize}%
\paragraph{Pontos de dúvida}%
\begin{itemize}%
\item%
\textit{(solution)} Item b, sinal da razão: o enunciado pede a razão entre E\_el e a 'variação da energia potencial gravitacional de P2'. Essa variação é negativa, porque P2 desce. A resolução responde E\_el/|$\Delta$U\_g| $\approx$ +1,6, embora tenha escrito $\Delta$U\_g = -2x\_s. O gabarito Poliedro dá -1,6. Se a banca exigir o sinal, o item pode perder nota. Convém ver com um professor se a resposta em módulo é aceita.%
\item%
\textit{(statement)} Item b, modelo quase estático: a mola está 'inicialmente relaxada', mas sobre P2 já atuam 2 N (peso) + 0,9 N (repulsão) para C. Se a carga for súbita, P2 oscila em torno do equilíbrio, e a compressão máxima pode chegar perto do dobro da estática, o que mudaria a razão. A resolução e o gabarito usam o equilíbrio estático (kx = 58/9 N). O verbo 'estime' parece justificar essa escolha.%
\item%
\textit{(statement)} Item c: a expressão 'após a condição da Figura 2' é ambígua. A resolução e o gabarito entendem que o vértice é atingido na fase livre da Figura 3 (BC horizontal, sem mola) e que C é o primeiro vértice. Com a mola presa, P2 não chegaria a C, porque k é muito alto.%
\item%
\textit{(key)} Na resolução do gabarito, a frase 'Uma constante que aparece o tempo todo:' vem seguida só de figuras, sem fórmula. Pelo contexto, é mg$\cdot$sen30$^\circ$ = 2 N. Isso não afeta as respostas.%
\end{itemize}%
\section*{Questão 8}%
\href{https://poliedroresolve.sistemapoliedro.com.br/vestibulares/ime/2027/ime-2-fase-dia-2-27-10-2026-2-fase-dia-2/dissertativa-8-fisica-2-fase-dia-2-ime-2027}{Resolução do gabarito}\par\smallskip%
\begin{longtable}{@{}p{0.06\textwidth} p{0.33\textwidth} p{0.33\textwidth} p{0.1\textwidth} p{0.08\textwidth}@{}}\toprule Item & Resolução & Gabarito & Veredito & Certo \\ \midrule
a & Afastando{-}se da rede: a contagem de máximos cai (7$\to$5$\to$3), o que só acontece quando D aumenta, porque y\_m = D$\cdot$tg$\theta$\_m cresce e as ordens mais altas saem pelas bordas. & O anteparo está se afastando da rede & igual & ambos \\
b & D(t1) = 0,2$\surd$7 $\approx$ 0,53 m & D1 = 0,53 m & igual & ambos \\
c & v = (0,3$\surd$3 - 0,1$\surd$7) m/s $\approx$ 0,25 m/s (0,255 m/s com raízes exatas; 0,254 m/s com a tabela) & v = 0,254 m/s & equivalente & ambos \\
d & $\Delta$s = 0,6$\surd$15 - 0,2$\surd$7 $\approx$ 1,79 m ($\Delta$t $\approx$ 7 s) & $\Delta$s = 1,792 m & equivalente & ambos \\
\bottomrule\end{longtable}%
\begin{itemize}%
\item%
\textbf{a)} Mesma conclusão e mesma justificativa. Com $\theta$\_m fixo, |y\_m| cresce com D, então as ordens saem do anteparo à medida que ele se afasta.%
\item%
\textbf{b)} Pela equação da rede, sen $\theta$3 = 3/4 e tg $\theta$3 = 3/$\surd$7. Impondo D$\cdot$3/$\surd$7 = 0,6 obtém{-}se D = 0,2$\surd$7 = 0,529 m (0,530 m com a tabela). Os dois valores coincidem.%
\item%
\textbf{c)} D2 = 0,6$\surd$3 = 1,039 m. Assim, v = (1,0392 - 0,5292)/2 = 0,2550 m/s com raízes exatas. Com a tabela, (1,038 - 0,530)/2 = 0,254 m/s. A diferença vem só do arredondamento das raízes, e a própria resolução cita 0,254 m/s com os valores da tabela.%
\item%
\textbf{d)} D1 = 0,6$\surd$15 = 2,324 m (2,322 m com a tabela). Então $\Delta$s = 2,3238 - 0,5292 = 1,795 m com raízes exatas, ou 2,322 - 0,530 = 1,792 m com a tabela. O valor 1,79 m da resolução bate com o do gabarito. O $\Delta$t $\approx$ 7 s é informação extra e está correto (1,7946/0,2550 = 7,04 s).%
\end{itemize}%
\paragraph{Pontos de dúvida}%
\begin{itemize}%
\item%
\textit{(solution)} Nas verificações, a resolução diz que a velocidade é de "centésimos de m/s", mas v $\approx$ 0,25 m/s está na casa dos décimos. É só um deslize de redação, sem efeito nos resultados.%
\item%
\textit{(key)} O gabarito usa as raízes arredondadas da tabela ($\surd$3 = 1,73, $\surd$7 = 2,65, $\surd$15 = 3,87) e chega a 0,254 m/s e 1,792 m. Com raízes exatas os valores são 0,255 m/s e 1,795 m. A diferença é só de arredondamento, mas vale ver se a banca exige os valores obtidos com a tabela.%
\item%
\textit{(statement)} Hipóteses implícitas que a resolução e o gabarito adotam: a ordem m = 4 ($\theta$ = 90$^\circ$, rasante) nunca chega ao anteparo, então os 7 máximos iniciais são m = 0, $\pm$1, $\pm$2, $\pm$3. Além disso, a franja na borda (y = L/2) conta como o instante exato da transição, e o anteparo é plano, com y = D$\cdot$tg$\theta$ (sem aproximação de pequenos ângulos).%
\end{itemize}%
\section*{Questão 9}%
\href{https://poliedroresolve.sistemapoliedro.com.br/vestibulares/ime/2027/ime-2-fase-dia-2-27-10-2026-2-fase-dia-2/dissertativa-9-fisica-2-fase-dia-2-ime-2027}{Resolução do gabarito}\par\smallskip%
\begin{longtable}{@{}p{0.06\textwidth} p{0.33\textwidth} p{0.33\textwidth} p{0.1\textwidth} p{0.08\textwidth}@{}}\toprule Item & Resolução & Gabarito & Veredito & Certo \\ \midrule
a & u = c($\surd$3$\cdot$h - L)/($\surd$3$\cdot$L - h) & u = c($\surd$3$\cdot$h - L)/($\surd$3$\cdot$L - h) & igual & ambos \\
b & L/$\surd$3 < h < L & L/$\surd$3 < h < L & igual & ambos \\
\bottomrule\end{longtable}%
\begin{itemize}%
\item%
\textbf{a)} Mesma expressão. Refiz a conta: a reta x = $\surd$3y faz 30$^\circ$ com o eixo x, e o raio refletido aponta em -y, então o desvio é de 120$^\circ$ e i = r = 30$^\circ$. A normal fica na direção (1/2, $\surd$3/2), a 30$^\circ$ da vertical, de modo que C = P + R$\cdot$(1/2, $\surd$3/2), com P = ct$\cdot$($\surd$3/2, 1/2). Daí saem ct$\cdot$$\surd$3/2 + ut/2 = L e ct/2 + ut$\cdot$$\surd$3/2 = h, que dão ct = $\surd$3L - h e ut = $\surd$3h - L. O sistema é o mesmo do gabarito. Conferi numericamente com (L,h) = (3,2), (1;0,7), (5;4,9) e (2;1,2): a menor raiz de |ctê - C| = ut é sempre $\surd$3L - h, ou seja, o primeiro contato é a entrada, e a direção refletida dá exatamente (0, -1).%
\item%
\textbf{b)} Mesma condição. As três exigências são ut > 0 $\Rightarrow$ h > L/$\surd$3; ct > 0 $\Rightarrow$ h < $\surd$3L, que já fica implícita; e u < c $\Rightarrow$ ($\surd$3+1)h < ($\surd$3+1)L $\Rightarrow$ h < L. A resolução vai além do gabarito: mostra que t1 é a menor raiz (f'(t1) < 0), logo é a entrada e não a saída, e descarta o outro sinal da normal. Isso só fortalece o resultado.%
\end{itemize}%
\paragraph{Pontos de dúvida}%
\begin{itemize}%
\item%
\textit{(statement)} O que conta como 'imposição física' no item b é ambíguo. Se a superfície se expande mais rápido que a componente normal da velocidade do fóton refletido (u > ($\surd$3/2)c, isto é, h > 5L/(3$\surd$3)), a esfera alcançaria o fóton logo depois da reflexão, já que a direção refletida foi calculada desprezando o movimento da superfície. Verifiquei: para $\tau$ pequeno, a diferença |r-C|$^2$ - R$^2$ $\approx$ ($\surd$3c - 2u)R1$\cdot$$\tau$. O enunciado manda desprezar esse efeito e lista apenas u < c como restrição. Por isso a resposta esperada (Poliedro) é L/$\surd$3 < h < L, e a resolução também a mantém, citando essa restrição mais forte só como observação. Vale confirmar com o professor se alguma banca aceitaria ou cobraria esse refinamento.%
\item%
\textit{(key)} O gabarito toma sem comentar a normal externa apontando para baixo e à esquerda (-(1/2, $\surd$3/2)) e não verifica que o ponto calculado é o primeiro contato (entrada), e não a saída. A resolução faz essa verificação. Isso não muda a resposta.%
\end{itemize}%
\section*{Questão 10}%
\href{https://poliedroresolve.sistemapoliedro.com.br/vestibulares/ime/2027/ime-2-fase-dia-2-27-10-2026-2-fase-dia-2/dissertativa-10-fisica-2-fase-dia-2-ime-2027}{Resolução do gabarito}\par\smallskip%
\begin{longtable}{@{}p{0.06\textwidth} p{0.33\textwidth} p{0.33\textwidth} p{0.1\textwidth} p{0.08\textwidth}@{}}\toprule Item & Resolução & Gabarito & Veredito & Certo \\ \midrule
a & Somente a Situação 2 (correntes iguais e no mesmo sentido, partícula descendo). Nas Situações 1 e 4, B no plano médio é diferente de zero e a força horizontal tem sinal fixo. Na Situação 3, a partícula sobe e o equilíbrio é instável. & Somente a situação 2 & igual & ambos \\
b & p $\approx$ ($\pi$d/2)$\cdot$$\surd$($\pi$mv/($\mu$0 i q)) = $\pi$v/$\omega$, com $\omega$ = (2/d)$\cdot$$\surd$($\mu$0 i q v/($\pi$m)) & T = $\pi$d$\cdot$$\surd$($\pi$m/($\mu$0 i q v)); p = ($\pi$d/2)$\cdot$$\surd$($\pi$mv/($\mu$0 i q)) & igual & ambos \\
\bottomrule\end{longtable}%
\begin{itemize}%
\item%
\textbf{a)} Conferi a resolução. Com correntes opostas (Situações 1 e 4), o campo no plano médio é B\_z(0) = -2$\mu$0 i/($\pi$d), diferente de zero. Assim já existe força horizontal no eixo, sempre no mesmo sentido, e não há restauração. Com correntes iguais (Situações 2 e 3), o campo é B\_z $\approx$ (4$\mu$0 i/($\pi$d$^2$))$\cdot$x e a força horizontal é F\_x = q$\cdot$?$\cdot$B\_z. Na Situação 2, a partícula desce (? = -v), então F\_x = -k$\cdot$x e a força é restauradora. Na Situação 3, ela sobe, então F\_x = +k$\cdot$x e a partícula se afasta do eixo. A justificativa física também confere: uma carga positiva descendo equivale a uma corrente para baixo, antiparalela às dos fios, e por isso é repelida pelo fio mais próximo. A figura do PDF confere com a transcrição. A resolução justifica cada uma das quatro situações, como o enunciado pede.%
\item%
\textbf{b)} Refiz a conta. A constante elástica efetiva é k = 4$\mu$0 i q v/($\pi$d$^2$), logo $\omega$$^2$ = 4$\mu$0 i q v/($\pi$md$^2$) e T = 2$\pi$/$\omega$ = $\pi$d$\cdot$$\surd$($\pi$m/($\mu$0 i q v)). O passo é a distância entre dois cruzamentos sucessivos do eixo y, ou seja, meio período: p = v$\cdot$T/2 = ($\pi$d/2)$\cdot$$\surd$($\pi$mv/($\mu$0 i q)). Essa é exatamente a expressão do gabarito. Também conferi numericamente (RK4 com o campo exato dos dois fios, m/q = 1e{-}8 kg/C, i = 10 A, d = 1 cm, v = 1 m/s, x0 = 10 $\mu$m): o espaçamento médio entre cruzamentos foi 0,78516 mm, contra 0,78540 mm da fórmula. O desvio relativo de -3,0$\cdot$10$^{-}$$^{4}$ coincide com a correção de segunda ordem -(3/16)($\omega$A/v)$^2$ que a resolução deduziu e que também verifiquei analiticamente. O T do gabarito confere com 2$\pi$/$\omega$.%
\end{itemize}%
\paragraph{Pontos de dúvida}%
\begin{itemize}%
\item%
\textit{(statement)} O enunciado não fornece g. A resolução e o gabarito desprezam o peso, o que é razoável, mas é uma hipótese implícita. A resolução a declara e estima o efeito (g/($\omega$v) \textasciitilde{} 10$^{-}$$^3$ no exemplo numérico).%
\item%
\textit{(statement)} Na figura, o passo p corresponde a meio período da oscilação em x (dois cruzamentos sucessivos do eixo y). Resolução e gabarito adotam essa leitura (p = vT/2). Quem tomasse p como o período completo (vT) teria resposta com o dobro do valor.%
\item%
\textit{(solution)} A resolução traz material além do pedido: integral primeira exata, órbitas das Situações 1 e 3, correção de segunda ordem e integração numérica. Conferi e está tudo correto. Só a correção -(3/16)($\omega$A/v)$^2$ ignora os termos de ordem (A/a)$^2$, que são desprezíveis quando $\kappa$ = 4$\pi$mv/($\mu$0 i q) ? 1, como no exemplo. Nada disso afeta a resposta final.%
\end{itemize}%
\end{document}